\section{Governance Cadence: Reviews Without Theatre}
Governance exists to make decisions, not to generate status theatre. Keep the cadence predictable and evidence-based.

\subsection{Governance cadence}
- Weekly portfolio review: top risks, blocked dependencies, key metrics.  
- Monthly product review: outcome progress, OKR status, learnings.  
- Quarterly planning: reset priorities, capacity, and major bets.

\subsection{Worked scenario: weekly portfolio review}
The portfolio review begins with blocked dependencies and critical risks. The TPM updates dependency owners. The PM adjusts Initiative status with evidence. The EM highlights capacity constraints. Decisions are logged as changes in Jira, not as notes in a deck.

\subsection{Case study insert}
In the weekly portfolio review, TS-117 is flagged as a blocker for EPIC-5. The TPM assigns an owner and the Jira link is updated during the meeting.
\subsection{Failure modes and fixes}
\begin{enumerate}
  \item \textbf{Failure:} Reviews become status recitals. \textbf{Fix:} Require decisions or escalations in every review.
  \item \textbf{Failure:} Metrics are too broad to act on. \textbf{Fix:} Focus on 3--5 metrics tied to outcomes.
  \item \textbf{Failure:} Quarterly plans ignore execution data. \textbf{Fix:} Bring throughput and cycle time into planning.
  \item \textbf{Failure:} Reviews are cancelled when busy. \textbf{Fix:} Treat reviews as the decision forum, not optional.
  \item \textbf{Failure:} Jira is updated after the meeting. \textbf{Fix:} Update issues during the review.
\end{enumerate}










